\chapter{اللقاء الثالث والثلاثون: نفقة المرضع وعدة الوفاة والإحداد}

\section{مقدمة}
السلام عليكم ورحمة الله وبركاته. الحمد لله رب العالمين، وأشهد أن لا إله إلا الله وحده لا شريك له، وأشهد أن محمداً عبده ورسوله، صلى الله عليه وعلى آله وصحبه وسلم.

استأنف الشيخ هذا اللقاء بعد انقطاع، من قوله تعالى: \begin{ayah}وَعَلَى الْمَوْلُودِ لَهُ رِزْقُهُنَّ وَكِسْوَتُهُنَّ بِالْمَعْرُوفِ\end{ayah}، ثم مضى في آيات الرضاع، وتعرض في أثناء ذلك لفائدة مهمة في باب الاستطاعة والقدر، ثم انتقل إلى أحكام عدة الوفاة والإحداد وما يتصل بها من لزوم البيت وترك الزينة.

ويمتاز هذا المجلس بأنه جمع بين الفقه الدقيق، والتحرير العقدي، وحسن جمع النصوص في الباب الواحد، وهي خصال ظاهرة جداً في صنيع \rawi{الطبري}.

\separator

\section{تأويل (وعلى المولود له رزقهن وكسوتهن بالمعروف)}
\isnad{قال أبو جعفر:} بيّن الله في هذه الآية أن نفقة المرضع وكسوتها على الأب، إذا كانت ترضع ولده، وأن ذلك يكون \textit{بالمعروف}، أي: على الوجه اللائق بحال الزوجين، من غير ظلم ولا تعسف ولا تكليف بما لا يحتمل.

وشرح \rawi{الطبري} أن \textit{المعروف} هنا مرده إلى ما يجب لمثلها على مثلِه، لأن الناس يتفاوتون في السعة والإقتار، وقد راعت الشريعة هذا التفاوت، فلم تجعل النفقة قدراً جامداً في جميع الأحوال، بل ربطتها بالوسع والطاقة.

ومن هنا ربط الشيخ هذه الآية بقوله تعالى: \begin{ayah}لِيُنْفِقْ ذُو سَعَةٍ مِنْ سَعَتِهِ وَمَنْ قُدِرَ عَلَيْهِ رِزْقُهُ فَلْيُنْفِقْ مِمَّا آتَاهُ اللَّهُ\end{ayah}، ليظهر أن أحكام النفقة مبنية على العدل والقدرة، لا على التماثل الصوري بين الناس.

\begin{fawaid}[رد الأحكام إلى القدرة والوسع من محاسن الشريعة]
لا تأتي الشريعة في أبواب النفقة والحقوق المالية بتكليف أعمى لا يراعي أحوال الناس، بل تبني الأمر على المعروف والوسع والطاقة، فيؤمر كل واحد بما يناسب حاله من غير ظلم ولا حيف.
\end{fawaid}

\section{تأويل (لا تكلف نفس إلا وسعها)}
وقف \rawi{الطبري} عند هذه الجملة وقفة مهمة، ففسرها بأن الله لا يحمل نفساً إلا ما لا يضيق عليها، ولا يتعذر عليها إذا أرادته. ثم رد على فهم باطل أدخله بعض أهل الكلام في هذا الباب، فجعلوا الوسع هنا دليلاً على معنى فاسد في مسألة القدرة والاستطاعة.

وبيّن الشيخ أن الاستطاعة في النصوص تأتي على معنيين:
\begin{itemize}
    \item استطاعة بمعنى الأهلية والقدرة قبل الفعل.
    \item واستطاعة بمعنى وقوع الاستجابة والفعل نفسه.
\end{itemize}

فقد يكون الإنسان مستطيعاً من جهة القدرة، ثم لا يفعل، فلا يلزم من ثبوت الاستطاعة الأولى حصول الثانية. وبذلك يزول كثير من الغلط في فهم نصوص التكليف ونصوص الذم للكافرين والمعرضين.

\begin{fawaid}[من المهم التمييز بين القدرة على الفعل ووقوع الفعل]
كثير من الغلط في باب القدر سببه الخلط بين استطاعة الأهلية، واستطاعة الامتثال والفعل. فإذا ميز القارئ بينهما، استقام له فهم نصوص كثيرة في التكليف والذم والوعد والوعيد.
\end{fawaid}

\section{فائدة عقدية: التكليف لا يتعلق إلا بالمقدور}
خرج الشيخ من هذا الموضع بقاعدة عظيمة: أن أمر الله ونهيه إنما يتعلقان بما للعبد سبيل إلى فعله أو تركه. أما ما لا يملكه من جهة الأصل، فليس محلاً للتكليف على هذا الوجه.

ولهذا لما انتقل \rawi{الطبري} بعد ذلك إلى الكلام على الجمع بين آية الرضاع وآية الحمل والفصال، قرر أن قوله تعالى: \begin{ayah}وَحَمْلُهُ وَفِصَالُهُ ثَلَاثُونَ شَهْرًا\end{ayah} ليس من باب الحد التعبدي الذي يطالب الخلق بتحصيله، بل من باب الخبر عن واقع من أحوال الناس. أما قوله: \begin{ayah}وَالْوَالِدَاتُ يُرْضِعْنَ أَوْلَادَهُنَّ حَوْلَيْنِ كَامِلَيْنِ\end{ayah} فهو في مقام التشريع والتقدير لمن أراد تمام الرضاعة.

\begin{fawaid}[ليس كل خبر في القرآن حكماً تكليفياً]
من الخطأ أن يحمل القارئ كل خبر في القرآن على أنه أمر أو حد تعبدي لازم. فبعض الأخبار يراد بها الإخبار عن الواقع أو عن بعض الصور، وبعضها يراد به التشريع والتحديد، ولا بد من التمييز بين المقامين.
\end{fawaid}

\section{تأويل (لا تضار والدة بولدها ولا مولود له بولده)}
وقف \rawi{الطبري} هنا عند القراءة والمعنى معاً، فذكر أن اللفظ يحتمل من جهة الصناعة وجهاً خبرياً ووجهاً نهياً، ثم رجح أن المقصود هو النهي: فلا يجوز أن تجعل الأم ولدها وسيلة إضرار بالأب، ولا أن يجعل الأب ولده وسيلة إضرار بالأم، بنزعٍ، أو تضييقٍ، أو امتناعٍ عن حق، أو إلزامٍ بما يقصد به المشقة لا المصلحة.

ومن حسن تقرير الشيخ في هذا الموضع أنه نبّه إلى أن النهي هنا متبادل، فلا يختص بجانب دون جانب. فالشريعة لا تنظر إلى مجرد ثبوت الحق، بل تمنع تحويل الحقوق نفسها إلى أبواب أذى وانتقام. ولذلك كان باب الرضاع - مع ما فيه من عاطفة الأمومة وحق الأبوة - من أشد الأبواب احتياجاً إلى التقوى والعدل.

\begin{fawaid}[قد يكون أصل الفعل مشروعاً ثم يحرم إذا صار وسيلة مضارة]
ليس كل ما ثبت فيه حق لأحد الطرفين يباح له استعماله كيف شاء، بل إذا نقل الحق من مقصوده الشرعي إلى الإضرار والانتقام انقلب ممنوعاً، وإن كان أصله مشروعاً.
\end{fawaid}

\section{تأويل (وعلى الوارث مثل ذلك)}
نقل \rawi{الطبري} في هذه الجملة خلافاً كثيراً: من الوارث؟ وما الذي يماثل المشار إليه بقوله: \textit{مثل ذلك}؟ ثم سلك طريقته المعروفة في تأخير الترجيح حتى يستوفي الاحتمالات.

وبيّن الشيخ أن هذا من المواضع التي يظهر فيها صبر \rawi{الطبري} على تفكيك الخلاف؛ فالنزاع ليس في كلمة واحدة فقط، بل في جهتين:
\begin{itemize}
    \item من المراد بالوارث: أهو وارث الصبي من جهة أبيه، أم ورثته عموماً، أم المولود نفسه؟
    \item وما المراد بالمماثلة: أهي النفقة والكسوة، أم منع المضارة، أم مجموع ما تقدم ذكره؟
\end{itemize}

ثم رجح \rawi{الطبري} أن المراد بالوارث هنا المولود نفسه، وأن الإشارة في \textit{مثل ذلك} إلى ما كان على الأب من مؤنة الرضاع إذا كان للمولود مال. ومهما يكن وجه الترجيح في المسألة، فالذي يهم هنا منهجياً أن \rawi{الطبري} أبطل وجوه التخصيص التي لا حجة عليها، فلم يرض أن يحمل اللفظ على بعض أفراده دون بعض إلا ببرهان.

\begin{fawaid}[منهج الطبري: لا تخصيص للفظ المحتمل بلا حجة]
إذا دار اللفظ العام أو المحتمل بين وجوه متعددة، لم يجز عند الطبري تقييده بواحد منها لمجرد الاستحسان، بل لا بد من دليل يخرجه من الاحتمال إلى التعيين.
\end{fawaid}

\separator

\section{تأويل (والذين يتوفون منكم ويذرون أزواجاً)}
انتقل \rawi{الطبري} بعد ذلك إلى عدة الوفاة، فقرر أن المتوفى عنها زوجها تتربص أربعة أشهر وعشراً، إلا أن تكون حاملاً فعدتها وضع الحمل. وهذا من المواضع التي تظهر فيها قوة جمع الآيات في الباب؛ إذ لا يستغني فهم آية البقرة هنا عن آية: \begin{ayah}وَأُولَاتُ الْأَحْمَالِ أَجَلُهُنَّ أَنْ يَضَعْنَ حَمْلَهُنَّ\end{ayah}.

وبيّن الشيخ أن التربص هنا ليس مجرد انتظار زمني مجرد، بل هو حكم تترتب عليه آثار معروفة في النكاح والزينة والمقام وغير ذلك مما جاءت به السنة.

\section{معنى التربص وحدوده}
فسر \rawi{الطبري} التربص في هذا الموضع بأنه احتباس المتوفى عنها بنفسها عن الأزواج، مع ما يلحق بذلك من أحكام الإحداد وترك الزينة، والبقاء في البيت الذي كانت فيه عند وفاة زوجها، على ما ثبتت به الأخبار.

وهنا نقل الخلاف في بعض هذه التفاصيل، فحكي قول من قصر التربص على ترك النكاح فقط، ثم رجح ما دلت عليه السنة من أن الإحداد وترك الزينة وترك النقلة عن المنزل داخل في معنى هذا الباب.

\begin{fawaid}[السنة تبين ما أجملته الآية]
قد يذكر القرآن الأصل الكلي، ثم تأتي السنة بتفصيل ما يدخل فيه من الصور والقيود. ولذلك لم يكتف الطبري هنا بظاهر لفظ التربص، بل رجع إلى السنة لبيان ما يندرج تحته من أحكام الإحداد والمقام.
\end{fawaid}

\section{الرد على من قصر العدة على ترك النكاح فقط}
ذكر \rawi{الطبري} قول من قال إن المتوفى عنها إنما نهيت عن الأزواج فقط، أما الطيب والزينة والخروج من منزلها فليست داخلة في الحكم. ثم رد ذلك بالأخبار الثابتة عن النبي صلى الله عليه وسلم في الإحداد، وبحديث \rawi{الفريعة بنت مالك} في لزوم البيت حتى يبلغ الكتاب أجله.

ومن محاسن هذا الموضع أن \rawi{الطبري} لم يترك الخلاف معلقاً، بل أبرز جهة الترجيح بوضوح:
\begin{itemize}
    \item ظاهر التنزيل يحتمل هذا المعنى العام.
    \item ثم جاءت السنة فبيّنت أن من التربص المشروع ترك الزينة والطيب، ولزوم البيت.
\end{itemize}

فصار العمل بالآية على وجهها إنما يتم بضم السنة إليها، لا بإفراد أحدهما عن الآخر.

\begin{fawaid}[الترجيح قد يكون برد قول الصحابي إلى الخبر المرفوع]
من قوة الطبري أنه لا يتوقف في رد اجتهاد من دونه أو فوقه إذا ثبت عنده الخبر عن النبي صلى الله عليه وسلم بخلافه. فالسنة عنده هي الحاكمة على ما سواها، ولو كان المخالف صحابياً.
\end{fawaid}

\section{حديث أسماء بنت عميس وكيف جمع الطبري بين الأخبار}
عرض \rawi{الطبري} ما قد يظن منه معارضته لأحاديث الإحداد، مثل الخبر الوارد في \rawi{أسماء بنت عميس}، ثم بيّن أن هذا لا يرفع أصل الإحداد، وإنما يحمل على ما يباح من اللباس الذي ليس زينة، وبذلك تجتمع الأخبار ولا تتدافع.

وهذا مثال جميل على طريقته في الجمع قبل دعوى التعارض: فإذا أمكن حمل كل خبر على وجه صحيح يوافق سائر النصوص، كان ذلك أولى من طرح بعضها أو إهدار دلالته.

\begin{fawaid}[الجمع بين النصوص أولى من دعوى التعارض]
إذا أمكن حمل الأخبار على وجوه متوافقة، لم يجز التسرع إلى التنازع بينها. وهذا من أعظم ما يحفظ الفقيه والمفسر من الاضطراب في أبواب الأحكام.
\end{fawaid}

\section{تأويل (أربعة أشهر وعشراً)}
ختم \rawi{الطبري} هذا الموضع بالتنبيه على أن العشرة هنا بالأيام مع لياليها، وأن هذا من الاستعمال العربي المعروف، فليس المراد مجرد عشرة ليال منفصلة عن أيامها. وفي هذا لمحة أخرى من لمحات التفسير: رد بعض تفاصيل الأحكام إلى لسان العرب واستعمالهم.

\section{تأويل (ولا جناح عليكم فيما عرضتم به من خطبة النساء)}
انتقل \rawi{الطبري} بعد ذلك إلى باب الخطبة في العدة، فقرر أن التعريض للمعتدة من وفاة زوجها جائز، بخلاف التصريح بعقد النكاح أو المواعدة المحرمة. وفسر التعريض بأنه من لحن الكلام الذي يفهم منه المراد من غير تصريح، كأن يظهر الرغبة في النكاح أو يثني عليها ثناءً يفهم منه القصد من غير أن يواجهها بعقد أو وعد صريح.

وهنا استخرج \rawi{الطبري} فائدة أصولية ولغوية لطيفة: أن التعريض غير التصريح حكماً، كما أنه غيره لفظاً. ولذلك لم يسوِّ بينهما في هذا الباب، واستنبط من الآية أن الأحكام قد تختلف باختلاف درجة البيان في الكلام، فلا يلحق المعرَّض فيه بالمصرَّح به دائماً.

\begin{fawaid}[اختلاف مراتب البيان يترتب عليه اختلاف الأحكام]
ليس كل ما فهم من الكلام يأخذ حكم ما صرح به؛ فبين التعريض والتصريح فروق معتبرة في اللسان والفقه، ومن لم يراعها اضطرب في تنزيل الأحكام.
\end{fawaid}

\section{تأويل (ولكن لا تواعدوهن سراً ولا تعزموا عقدة النكاح)}
ذكر \rawi{الطبري} الخلاف في معنى \textit{السر} هنا، ثم رجح أن المراد به الجماع وما يتصل بالمواعدة الخفية عليه، لا مجرد الكتمان في القول. ووجه ذلك بأن العرب تسمي الغشيان سراً لوقوعه في الخفاء، ولأن حمل الآية على مجرد المعاهدة الخفية يوقع في اضطراب المعنى، إذ يصير الممنوع هو الكتمان فقط، مع أن أصل المواعدة نفسها على هذا الوجه ممنوع سراً وعلانية.

ثم فسر قوله تعالى: \begin{ayah}وَلَا تَعْزِمُوا عُقْدَةَ النِّكَاحِ\end{ayah} بمعنى: لا تصححوا عقد النكاح، ولا توقعوه، ولا تمضوه حتى يبلغ الكتاب أجله. فالمباح في هذه المرحلة هو التعريض المأذون، لا إنشاء العقد، ولا قطع المواعدة التي تتجاوز حد الإباحة.

\section{تأويل (واعلموا أن الله يعلم ما في أنفسكم فاحذروه)}
ختم \rawi{الطبري} هذا الموضع بما يردّ الأحكام كلها إلى أصلها الإيماني: أن الله يعلم ما في النفوس من الميل والهوى، ومن إرادة النكاح، ومن مجاوزة الحد أو الوقوف عنده. ولذلك جاء الأمر بالحذر هنا؛ لأن كثيراً من هذا الباب يقع في المساحات الخفية التي قد يهون أمرها على الناس، لكنها لا تخفى على الله.

ثم أتبع ذلك بقوله تعالى: \begin{ayah}وَاعْلَمُوا أَنَّ اللَّهَ غَفُورٌ حَلِيمٌ\end{ayah}، ليجمع بين التحذير والرجاء. فليس المقصود مجرد التخويف، بل أن يردع المكلف عن المخالفة، ويُفتح له مع ذلك باب التوبة إن زل أو قارب ما لا ينبغي.

\separator

خلاصة هذا اللقاء أن \rawi{الطبري} جمع فيه بين فقه النفقة، والتحرير العقدي في باب الاستطاعة، ثم أحكام عدة الوفاة والإحداد، مع تقرير أصل منهجي ثابت: أن القرآن لا يفهم في هذه الأبواب إلا بجمع آياته، وضم السنة الصحيحة إليها، والنظر في لسان العرب، ورد المتشابه إلى المحكم.
